\input{noncyclic/cusp_image_3_15_D4.tex}
\input{noncyclic/exact_3_15_D4_specialization.tex}
\begin{proposition}\label{prop:octic-315-complete}
The group $C_3\times C_{15}$ occurs as the full torsion subgroup
of an elliptic curve over an octic field. Every elliptic curve over
an octic field containing this group has full torsion $C_3\times C_{15}$.
\end{proposition}
\begin{proof}
Let $z^2+z+1=0$ and put $F=\mathbb Q(z)$.
The polynomial
\[
 f(t)=t^4-5t^3+8t^2-6t+3
\]
is irreducible over $F$, so $K=F(w)$, with $f(w)=0$, is octic.
Consider the elliptic curve
\[
 E:y^2+a_1xy+a_3y=x^3,
\]
where
\[
\begin{split}
 a_1={}&\frac{(-7z-8)w^3+(23z+34)w^2+(-8z-43)w}{6}
                         +\frac{-3z+10}{2},\\
 a_3={}&\frac{(17z+13)w^3+(-82z-71)w^2+(106z+119)w}{12}
                         -2z-5.
\end{split}
\]
Then $P=(0,0)$ has order $3$. The point $Q=(x_Q,y_Q)$, with
\[
\begin{split}
 x_Q={}&\frac{(-z+1)w^3+(8z-2)w^2+(-17z-1)w}{6}
                              +\frac{3z}{2},\\
 y_Q={}&\frac{(z+2)w^3+(-4z-8)w^2+(3z+6)w+2z+3}{4},
\end{split}
\]
has exact order $15$, and $e_3(P,5Q)=z^2$.
These assertions are exact identities in $K$;
the discriminant of $E$ is nonzero and $j(E)=1331/8$.
Thus $E(K)$ contains $C_3\times C_{15}$.

The discriminant of $f$ is $-507$, so the prime of $K$ specified
by $z\mapsto6$, $w\mapsto4$ above $43$ is unramified, with
residue field $\mathbb F_{43}$. The curve has good reduction
there, with group $[3,15]$ of order $45$. Good reduction over
the unramified odd local field injects on all torsion.
Consequently $E(K)_{\mathrm{tors}}=C_3\times C_{15}$.
The independent certificate checks irreducibility over $F$,
the finite point count and the $45$ elements generated by the
reductions of $P,Q$, in addition to the native exact level check.

Put $X=X_1(3,15)$. The canonical isomorphism invariant $\tau=-y_Q^2/x_Q^3$ of
the triple $(E,P,Q)$ has degree $4$ over $F$: the coefficient
vectors of $1,\tau,\tau^2,\tau^3$ have rank $4$ over $F$.
Here $P$ is at the origin with tangent $y=0$, so the remaining
coordinate changes are scalings, which preserve $\tau$. Thus the displayed fine
moduli point has residue field exactly $K$. Its trace is a
closed $F$-rational effective divisor $D$ of degree $4$ on $X$.
Lemma~\ref{lem:octic-315-exact-specialization} puts the class
of $D$ in a single orbit of twelve degree-$4$ classes.
Now let an arbitrary octic elliptic curve contain $C_3\times C_{15}$.
Its field contains $F$ by the Weil pairing. Taking the trace of
its moduli point, with multiplicity if necessary, gives another
noncuspidal effective $F$-rational degree-$4$ divisor $D'$.
By that lemma there is an $F$-defined automorphism $g$ with
$[D']=[gD]$. Since $\operatorname{gon}_F(X)\geq6>4$, an
$F$-rational divisor class of degree $4$ has at most one effective
representative. Hence $D'=gD$.

The $48$ automorphisms used in the lemma change only the level
structure on the same elliptic curve. The universal curve at
$gD$ is therefore the curve just exhibited, up to field
conjugacy and isomorphism. The fine moduli interpretation
identifies these curves over their residue fields. In particular
the full torsion is again $C_3\times C_{15}$.
\end{proof}

\begin{corollary}\label{cor:octic-no330}
No elliptic curve over an octic field contains $C_3\times C_{30}$.
\end{corollary}
\begin{proof}
Such a group strictly contains $C_3\times C_{15}$, contradicting
Proposition~\ref{prop:octic-315-complete}.
\end{proof}
